\begin{explanationbox}
	\textbf{\textcolor{CExplanation}{一句话直觉}}
	用离心角连接坐标，使椭圆问题转化为三角运算。

	\medskip
	\textbf{\textcolor{CExplanation}{核心拆解}}
	\begin{description}[style=nextline, leftmargin=2.8em, labelsep=0.5em, itemsep=0.45em]
		\item[\textbf{要点一（角度参数化）：}]
			设 $x=a\cos\theta,\;y=b\sin\theta$，统一用 $\theta$ 表示坐标。
		\item[\textbf{要点二（三角恒等式）：}]
			代入椭圆方程恒成立，因 $\cos^2\theta+\sin^2\theta=1$。
		\item[\textbf{要点三（参数范围）：}]
			$\theta\in\mathbb{R}$ 可覆盖整个椭圆，通常取 $[0,2\pi)$。
	\end{description}

	\medskip
	\textbf{\textcolor{CExplanation}{几何本质（辅助圆模型）}}
	椭圆可视为半径为 $a$ 的大圆沿 $y$ 轴压缩 $\frac{b}{a}$ 倍，$\theta$ 是大圆上的圆心角。

	\medskip
	\textbf{\textcolor{CExplanation}{代数意义（平方和恒等）}}
	参数方程直接满足 $\frac{x^2}{a^2}+\frac{y^2}{b^2}=1$，是三角恒等式的几何体现。

	\medskip
	\textbf{\textcolor{CExplanation}{双视角理解（椭圆上点到中心的距离）}}
	设椭圆 $\frac{x^2}{a^2}+\frac{y^2}{b^2}=1\;(a>b>0)$，中心为固定点 $O(0,0)$，$P$ 为椭圆上的动点。研究 $OP$ 的最大值和最小值，可以从两个相反方向观察。
	\begin{description}[style=nextline, leftmargin=2.8em, labelsep=0.5em, itemsep=0.5em]
		\item[\textbf{视角一（由内到外：从不动点 $O$ 出发）：}]
			先固定 $O$，以 $O$ 为圆心画圆 $x^2+y^2=r^2$，再让半径 $r$ 从小到大变化。$r<b$ 时，圆还碰不到椭圆；$r=b$ 时，圆第一次在 $(0,\pm b)$ 与椭圆相切，故最小距离为 $b$；$b<r<a$ 时，圆与椭圆有四个交点；$r=a$ 时，圆在 $(\pm a,0)$ 与椭圆相切，达到最大距离 $a$；$r>a$ 时不再相交。
			\par\textbf{直觉：} 不让圆心动，只改变半径；第一次碰到与最后一次碰到椭圆的半径，就是距离的两个边界。
		\item[\textbf{视角二（由外到内：从动点 $P$ 出发）：}]
			先让 $P$ 在椭圆上运动，再从 $P$ 连线指向固定的 $O$，把每个位置的距离 $OP$ 写成函数。设 $P(a\cos\theta,b\sin\theta)$，则
			\[
				\begin{aligned}
					OP^2&=(a\cos\theta)^2+(b\sin\theta)^2\\
					&=b^2+(a^2-b^2)\cos^2\theta.
				\end{aligned}
			\]
			由于 $0\le\cos^2\theta\le1$，得到 $b\le OP\le a$。当 $P=(0,\pm b)$ 时取最小值，当 $P=(\pm a,0)$ 时取最大值。
			\par\textbf{直觉：} 动的是椭圆上的点，测量方向是从 $P$ 指向不动的 $O$；不要求 $P$ 在整个运动过程中离 $O$ 越来越近。
	\end{description}
	\textbf{两种视角的统一：} 第一种从固定圆心寻找与椭圆的临界相切圆，第二种从椭圆动点求距离函数的最值；观察方向相反，结果相同：
	\[
		\boxed{OP_{\min}=b,\qquad OP_{\max}=a.}
	\]
	\textbf{特别区分：} 半径为 $a$、$b$ 的两个同心圆分别与椭圆相切，并不是这两个圆互相相切；这里不断改变半径的\textbf{距离圆}，也不要与前面用于生成参数方程的\textbf{辅助圆}混为一谈。

	\medskip
	\textbf{\textcolor{CExplanation}{考点价值（求最值与轨迹）}}
	\begin{itemize}[leftmargin=*, itemsep=0.5em, topsep=0.4em]
		\item \textbf{考法一（最值问题）：} 设点为 $(a\cos\theta,b\sin\theta)$，利用三角有界性求最值。
		\item \textbf{考法二（参数范围）：} 由 $\theta$ 范围确定变量范围，或消参得普通方程。
	\end{itemize}

	\medskip
	\textbf{\textcolor{CExplanation}{顿悟点}}
	离心角 $\theta$ 不是极角 $\varphi$，关系为 $\tan\varphi=\frac{b}{a}\tan\theta$。

	\medskip
	\textbf{\textcolor{CExplanation}{使用场景}}
	\begin{itemize}[leftmargin=*, itemsep=0.5em, topsep=0.4em]
		\item \textbf{场景一（设点简化）：} 求椭圆上点到直线距离、面积最值时，用参数设点减少变量。
		\item \textbf{场景二（消参转化）：} 处理椭圆与直线综合问题时，可转化为三角函数问题。
	\end{itemize}

\end{explanationbox}
